\chapter{Bijou Vision}\label{record-035}

\noindent{\footnotesize\color{BookMuted}\path{docs/strategy/bijou-vision.md}\par}\smallskip

\begin{quote}
Historical note: the canonical project-identity signpost is now \href{https://github.com/flyingrobots/bijou/blob/cbd70612140e7a450d764b3c7be1c6af878b99bc/docs/VISION.md}{\texttt{d\allowbreak{}o\allowbreak{}c\allowbreak{}s\allowbreak{}/\allowbreak{}V\allowbreak{}I\allowbreak{}S\allowbreak{}I\allowbreak{}O\allowbreak{}N\allowbreak{}.\allowbreak{}m\allowbreak{}d}}. This file remains as expanded background and provenance, not as the active queue surface.
\end{quote}

\emph{Living statement of what Bijou is for, how it should feel, and where it should be used}

\section{What Bijou is}\label{record-035-what-bijou-is}

Bijou is a framework for building \textbf{surface-native terminal applications} that are:

\begin{itemize}
\tightlist
\item
  styled
\item
  stateful
\item
  testable
\item
  humane
\item
  honest across different output modes
\end{itemize}

Bijou is not trying to pretend the terminal is a browser.

It is trying to take the terminal seriously as its own medium:

\begin{itemize}
\tightlist
\item
  dense, but not cramped
\item
  expressive, but not decorative for its own sake
\item
  interactive, but still degradable
\item
  structured, inspectable, and testable
\end{itemize}

The long-term goal is not only to ship components.

The goal is to make it straightforward to build terminal software that feels:

\begin{itemize}
\tightlist
\item
  calm
\item
  explicit
\item
  reliable
\item
  accessible
\item
  governable
\end{itemize}

\section{What Bijou is not}\label{record-035-what-bijou-is-not}

Bijou is not primarily:

\begin{itemize}
\tightlist
\item
  a string-templating helper for ad hoc CLI output
\item
  a browser-emulation layer in a terminal window
\item
  a pile of isolated widgets without shell doctrine
\item
  a framework that treats docs, accessibility, and graceful degradation as afterthoughts
\end{itemize}

It should stay willing to say ``this belongs to the shell,'' ``this belongs to layout,'' or ``this should lower differently in a non-rich mode'' instead of flattening every concern into raw component calls.

\section{Product vision}\label{record-035-product-vision}

Bijou should become the framework people reach for when they want to build serious terminal software with:

\begin{itemize}
\tightlist
\item
  real application state
\item
  real interaction structure
\item
  real layout ownership
\item
  real testing and replayability
\item
  real design standards
\end{itemize}

The ideal Bijou app should be explainable in a few plain statements:

\begin{itemize}
\tightlist
\item
  where state lives
\item
  what views are visible
\item
  what currently owns input
\item
  how layout is being constrained
\item
  what commands can change state
\item
  what effects may happen as a consequence
\end{itemize}

That clarity is part of the product.

\section{Design philosophy}\label{record-035-design-philosophy}

\subsection{1. Hexagonal at the core}\label{record-035-hexagonal-at-the-core}

The core should own:

\begin{itemize}
\tightlist
\item
  state
\item
  view stack
\item
  layout
\item
  routing
\item
  commands
\item
  effects
\end{itemize}

Adapters should own:

\begin{itemize}
\tightlist
\item
  terminal protocols
\item
  Node boundaries
\item
  clocks, files, and platform APIs
\item
  transport and environment details
\end{itemize}

If a behavior is important enough to test, inspect, or lower accessibly, it should probably exist inside the hexagon.

\subsection{2. Graceful lowering is part of the contract}\label{record-035-graceful-lowering-is-part-of-the-contract}

Rich mode is not the only reality.

Bijou should preserve meaning across:

\begin{itemize}
\tightlist
\item
  rich interactive mode
\item
  static mode
\item
  pipe mode
\item
  accessible mode
\end{itemize}

Visual parity is not the requirement. Semantic honesty is.

\subsection{3. Layout is a first-class concern}\label{record-035-layout-is-a-first-class-concern}

Components should not only know how to render.

They should also have explicit answers for:

\begin{itemize}
\tightlist
\item
  how they size
\item
  how they stretch or shrink
\item
  how they align and anchor
\item
  how they overflow
\item
  when they need a viewport
\end{itemize}

If layout ownership is vague, interaction ownership becomes vague too.

\subsection{4. Input belongs to the active interaction model}\label{record-035-input-belongs-to-the-active-interaction-model}

Input should flow through explicit ownership:

\begin{itemize}
\tightlist
\item
  topmost active view/layer first
\item
  layout-defined interaction geometry
\item
  component-local handling
\item
  commands and effects emitted intentionally
\end{itemize}

This is how modals block lower layers, how shells avoid stale controls, and how input remains explainable instead of branch-order magic.

\subsection{5. Commands and effects are not the same thing}\label{record-035-commands-and-effects-are-not-the-same-thing}

Commands mean:

\begin{itemize}
\tightlist
\item
  ``app, change something meaningful''
\end{itemize}

Effects mean:

\begin{itemize}
\tightlist
\item
  ``runtime, perform this consequence''
\end{itemize}

Bijou should keep that distinction clean.

It is not just type hygiene. It is how stateful behavior stays understandable.

\subsection{6. The shell is a product surface, not glue}\label{record-035-the-shell-is-a-product-surface-not-glue}

Search, help, settings, notifications, quit flows, and framing chrome are part of the user experience.

They should not feel like bolted-on conveniences around ``the real app.''

\subsection{7. The docs are the proving ground}\label{record-035-the-docs-are-the-proving-ground}

DOGFOOD is important because it forces Bijou to live with its own decisions.

The docs should not only describe the system. They should exercise it.

\subsection{8. Human and agent users both matter}\label{record-035-human-and-agent-users-both-matter}

Agents are first-class users of the framework and its artifacts.

A good Bijou system should be understandable enough that:

\begin{itemize}
\tightlist
\item
  humans can reason about it
\item
  agents can inspect it
\item
  tests can prove it
\item
  docs can teach it
\end{itemize}

That means explicit contracts beat folklore.

\section{Design values}\label{record-035-design-values}

\subsection{Calm over noise}\label{record-035-calm-over-noise}

Density is good. Chaos is not.

\subsection{Explicit over magical}\label{record-035-explicit-over-magical}

The framework should prefer clear ownership and named seams over hidden cleverness.

\subsection{Honest over flashy}\label{record-035-honest-over-flashy}

A simpler truthful fallback is better than a rich presentation that lies about capability.

\subsection{Structured over stringly}\label{record-035-structured-over-stringly}

Structured surfaces, layout nodes, and runtime objects are preferable where behavior needs to be inspectable or composable.

\subsection{Testable by design}\label{record-035-testable-by-design}

Important behavior should be expressible in deterministic tests instead of only in screenshots or subjective demos.

\subsection{Humane by default}\label{record-035-humane-by-default}

The framework should make it easier to do the right thing for:

\begin{itemize}
\tightlist
\item
  accessibility
\item
  localization
\item
  explainability
\item
  reduced motion
\item
  low cognitive load
\end{itemize}

\section{Recommended use cases}\label{record-035-recommended-use-cases}

Bijou is a strong fit for:

\subsection{Operator and internal tools}\label{record-035-operator-and-internal-tools}

Apps where people live in terminals and need:

\begin{itemize}
\tightlist
\item
  dashboards
\item
  workflow tools
\item
  review surfaces
\item
  deployment or environment controls
\item
  data inspection
\item
  incident response interfaces
\end{itemize}

\subsection{Developer tools}\label{record-035-developer-tools}

Apps where state, logs, structured output, and keyboard-driven flows matter:

\begin{itemize}
\tightlist
\item
  local dev control planes
\item
  test/replay tools
\item
  code review helpers
\item
  repo/workspace tools
\item
  build and release interfaces
\end{itemize}

\subsection{Guided and explainable interfaces}\label{record-035-guided-and-explainable-interfaces}

Apps that need to teach, justify, or walk users through decisions:

\begin{itemize}
\tightlist
\item
  AI-assisted flows
\item
  operator handoff tools
\item
  approvals
\item
  audits
\item
  onboarding/tutorial surfaces
\end{itemize}

\subsection{Dense information browsing}\label{record-035-dense-information-browsing}

Apps that need to browse or compare structured information:

\begin{itemize}
\tightlist
\item
  tables
\item
  trees
\item
  timelines
\item
  layered detail panes
\item
  searchable docs and knowledge systems
\end{itemize}

\subsection{Documentation and field guides}\label{record-035-documentation-and-field-guides}

Apps where the documentation itself should be interactive, testable, and mode-aware.

DOGFOOD is the canonical example of this.

\section{Less ideal use cases}\label{record-035-less-ideal-use-cases}

Bijou is probably a weaker fit when the primary requirement is:

\begin{itemize}
\tightlist
\item
  pixel-perfect graphical presentation
\item
  highly freeform mouse-first visual design
\item
  document editing that depends on browser/editor ecosystems Bijou does not aim to replace
\item
  experiences where the terminal itself is already the wrong medium
\end{itemize}

That is not a failure. It is a boundary.

\section{What good Bijou apps should feel like}\label{record-035-what-good-bijou-apps-should-feel-like}

A good Bijou app should answer these questions quickly:

\begin{itemize}
\tightlist
\item
  Where am I?
\item
  What can I do right now?
\item
  What currently owns input?
\item
  What changes if I activate this?
\item
  How do I dismiss or back out safely?
\end{itemize}

It should also make these answers legible to the builder:

\begin{itemize}
\tightlist
\item
  Which state is active?
\item
  Which views are stacked?
\item
  Which layout owns this space?
\item
  Which component owns this interaction?
\item
  Which command changes state?
\item
  Which effect leaves the core?
\end{itemize}

\section{Relationship to other docs}\label{record-035-relationship-to-other-docs}

This document sits above more specific doctrine and design notes:

\begin{itemize}
\tightlist
\item
  \href{https://github.com/flyingrobots/bijou/blob/cbd70612140e7a450d764b3c7be1c6af878b99bc/docs/strategy/bijou-ux-doctrine.md}{Bijou UX Doctrine}
\item
  \href{https://github.com/flyingrobots/bijou/blob/cbd70612140e7a450d764b3c7be1c6af878b99bc/docs/strategy/humane-shell.md}{The Humane Shell}
\item
  \href{https://github.com/flyingrobots/bijou/blob/cbd70612140e7a450d764b3c7be1c6af878b99bc/docs/strategy/accessibility-and-assistive-modes.md}{Accessibility and Assistive Modes}
\item
  \href{https://github.com/flyingrobots/bijou/blob/cbd70612140e7a450d764b3c7be1c6af878b99bc/docs/strategy/localization-and-bidirectionality.md}{Localization and Bidirectionality}
\item
  \href{https://github.com/flyingrobots/bijou/blob/cbd70612140e7a450d764b3c7be1c6af878b99bc/docs/strategy/ai-explainability-standard.md}{AI Explainability Standard}
\item
  \href{https://github.com/flyingrobots/bijou/blob/cbd70612140e7a450d764b3c7be1c6af878b99bc/docs/strategy/content-guide.md}{Content Guide}
\item
  \href{https://github.com/flyingrobots/bijou/blob/cbd70612140e7a450d764b3c7be1c6af878b99bc/docs/ARCHITECTURE.md}{Architecture --- Bijou}
\end{itemize}

\section{Living-document note}\label{record-035-living-document-note}

This document should evolve when Bijou's self-conception changes:

\begin{itemize}
\tightlist
\item
  if the architecture center of gravity changes
\item
  if the product promise changes
\item
  if the recommended use cases change
\item
  if the framework starts aiming at a different class of software
\end{itemize}

If those things change and this file stays stale, the repo will start lying about what Bijou is for.

